%----------------------------------------------------------------------------------------
% Tailored CV — Machine Learning Scientist, Multimodal AI
% Natera (Remote). Greenhouse Job 6004385004
% Reusable template for multimodal-AI / ML-scientist (oncology) roles.
%----------------------------------------------------------------------------------------

\documentclass[10pt,a4paper,sans]{moderncv}

\usepackage[utf8]{inputenc}
\usepackage[T1]{fontenc}

\moderncvstyle{classic}
\moderncvcolor{blue}

\usepackage[scale=0.78]{geometry}
\usepackage{ragged2e}

%----------------------------------------------------------------------------------------

\firstname{Kundai Farai}
\familyname{Sachikonye}

\address{Auf der Lichtung 19}{82178 Puchheim, Germany}
\mobile{(+49) 1626386344}
\email{kundai.sachikonye@bitspark.com}
\homepage{github.com/fullscreen-triangle}{github.com/fullscreen-triangle}

\quote{Multimodal machine learning for oncology --- integrating imaging,
molecular, and clinical data, from PyTorch prototype to deployed, validated tool.}

%----------------------------------------------------------------------------------------

\begin{document}

\makecvtitle

%----------------------------------------------------------------------------------------
%	EXPERIENCE
%----------------------------------------------------------------------------------------

\section{Experience}

\cventry{2021--present}{Independent ML \& Computational Biology Researcher}{Bitspark GmbH}{}{}{
  Full-time development of multimodal machine-learning systems across genomics,
  proteomics/metabolomics, and biomedical imaging. PyTorch model development,
  distributed training pipelines, and production deployment. All systems publicly
  documented, validated against reference data, and deployed (browser-native /
  GPU-accelerated).
}

\cventry{2019--2021}{Computational Lipidomics --- ML \& HPC}{TUM / Bayerisches Forschungsinstitut}{Munich}{}{
  Deep-learning mass-spectrometry analysis at HPC scale: CNN-based spectral
  classification (PyTorch) alongside numerical pipelines; multi-omics statistical
  integration; distributed processing of $>$100\,GB datasets.
}

\cventry{2017--2018}{Glycomics and Proteomics}{Universitätsklinikum Hamburg-Eppendorf}{Hamburg}{}{
  Automated LC-MS/MS and MALDI-TOF data pipelines; reproducible workflows for
  high-throughput clinical instrument data.
}

\cventry{2013}{Cancer Biomarker Discovery}{University Medical Center Groningen}{Groningen}{}{
  UPLC-MS/MS shotgun proteomics on tumour samples; statistical analysis of
  high-dimensional feature sets.
}

%----------------------------------------------------------------------------------------
%	EDUCATION
%----------------------------------------------------------------------------------------

\section{Education}

\cventry{2019--2021}{Doctoral research, Computational Lipidomics}{Technische Universität München}{}{}{
  HPC mass-spectrometry ML pipelines, multi-omics modelling, distributed training.
  Transitioned to independent research.
}
\cventry{2018--2019}{MSc Computational Biology}{Universität Tübingen}{}{}{
  Sequence analysis, variant calling, RNA-seq, ML for biological data.
}
\cventry{2014--2017}{MSc Pharmaceutical Biotechnology}{HAW Hamburg}{}{}{}
\cventry{2011--2014}{BSc Biochemistry and Cell Biology}{Jacobs University Bremen}{}{}{}

%----------------------------------------------------------------------------------------
%	SELECTED SYSTEMS
%----------------------------------------------------------------------------------------

\section{Selected ML Systems}

\cvitem{lavoisier}{
  \textbf{Deep-Learning Mass-Spectrometry Pipeline.}
  Dual pipeline: numerical MS analysis + CNN spectral classification (PyTorch);
  Ray-distributed execution; memory-mapped I/O and HDF5 for $>$100\,GB datasets;
  98.9\% NIST feature-extraction accuracy. Production-grade, deployed.
  \href{https://github.com/fullscreen-triangle/lavoisier}{github.com/fullscreen-triangle/lavoisier}
}

\cvitem{gospel}{
  \textbf{Multi-Omics Integration.}
  Joint Bayesian-network and mixture modelling over host genotype, transcriptomics,
  and metagenomics; $10^6$ variants/s; validated against 1000~Genomes, gnomAD,
  UK~Biobank. Multimodal molecular integration at scale.
  \href{https://github.com/fullscreen-triangle/gospel}{github.com/fullscreen-triangle/gospel}
}

\cvitem{syndrome}{
  \textbf{Computational Oncology / Immunopeptidomics.}
  Neoantigen identification and MHC-binding prediction (AUC~0.81 on 2{,}847 IEDB
  peptides), tumour disease-state characterisation. Browser-native (WebGPU).
  \href{https://syndrome-seven.vercel.app/tools}{syndrome-seven.vercel.app}
}

%----------------------------------------------------------------------------------------
%	TECHNICAL SKILLS
%----------------------------------------------------------------------------------------

\section{Technical Skills}

\cvitem{ML / Deep Learning}{
  PyTorch (CNNs, sequence models, attention); foundation-model fine-tuning
  (LoRA adapters); knowledge distillation; scikit-learn, XGBoost; Bayesian-network
  and mixture modelling
}

\cvitem{Multimodal biomedical data}{
  Genomics (variant calling, RNA-seq); proteomics / metabolomics (mass spectrometry);
  microscopy image analysis (segmentation, deconvolution, morphometry ---
  digital-pathology-adjacent); immunopeptidomics; integration of imaging +
  molecular + clinical modalities
}

\cvitem{Cloud / Distributed}{
  Ray, Dask, SLURM-compatible pipelines; Docker, Kubernetes; cloud-portable
  training workflows; memory-mapped I/O, HDF5, Parquet
}

\cvitem{Languages}{
  Python (primary), Rust, R, Bash, SQL, TypeScript
}

\cvitem{Reference data}{
  1000~Genomes, gnomAD, UK~Biobank, IEDB, NIST
}

%----------------------------------------------------------------------------------------
%	LANGUAGES
%----------------------------------------------------------------------------------------

\section{Languages}

\cvitemwithcomment{English}{Native / Fluent}{}
\cvitemwithcomment{German}{Conversational}{resident in Germany since 2011}

%----------------------------------------------------------------------------------------

\renewcommand{\listitemsymbol}{-~}

\end{document}
